\documentclass[aspectratio=169,12pt]{beamer}
% Estilo canónico para las presentaciones de Métodos Numéricos 2026-2.
\usepackage[T1]{fontenc}
\usepackage[utf8]{inputenc}
\usepackage[spanish,es-nodecimaldot]{babel}
\usepackage{lmodern}
\usepackage{amsmath,amssymb,mathtools}
\usepackage{booktabs}
\usepackage{graphicx}
\usepackage{xcolor}
\usepackage{hyperref}

\definecolor{upazul}{HTML}{003B70}
\definecolor{upazulclaro}{HTML}{EAF2F8}
\definecolor{updorado}{HTML}{C5962E}
\definecolor{upgris}{HTML}{4A5560}
\definecolor{upverde}{HTML}{167C5A}
\definecolor{uprojo}{HTML}{B33A3A}

\setbeamercolor{normal text}{fg=black,bg=white}
\setbeamercolor{structure}{fg=upazul}
\setbeamercolor{frametitle}{fg=white,bg=upazul}
\setbeamercolor{title}{fg=upazul}
\setbeamercolor{subtitle}{fg=upgris}
\setbeamercolor{block title}{fg=white,bg=upazul}
\setbeamercolor{block body}{fg=black,bg=upazulclaro}
\setbeamercolor{alerted text}{fg=uprojo}
\setbeamercolor{example text}{fg=upverde}

\setbeamerfont{title}{size=\Huge,series=\bfseries}
\setbeamerfont{subtitle}{size=\Large}
\setbeamerfont{frametitle}{size=\LARGE,series=\bfseries}
\setbeamerfont{normal text}{size=\large}
\setbeamerfont{block title}{series=\bfseries}

\setbeamertemplate{navigation symbols}{}
\setbeamertemplate{itemize item}{\color{updorado}$\blacktriangleright$}
\setbeamertemplate{itemize subitem}{\color{upazul}$\bullet$}
\setbeamertemplate{blocks}[rounded][shadow=false]

\setbeamertemplate{frametitle}{%
  \nointerlineskip
  \begin{beamercolorbox}[wd=\paperwidth,ht=0.78cm,dp=0.24cm,leftskip=0.55cm]{frametitle}
    \usebeamerfont{frametitle}\insertframetitle
  \end{beamercolorbox}}

\setbeamertemplate{footline}{%
  \leavevmode
  \hbox{%
    \begin{beamercolorbox}[wd=.82\paperwidth,ht=2.7ex,dp=1.2ex,leftskip=.5cm]{author in head/foot}
      \color{upgris}\scriptsize Métodos Numéricos 2026-2 \hspace{0.8em} | \hspace{0.8em} Universidad del Pacífico
    \end{beamercolorbox}%
    \begin{beamercolorbox}[wd=.18\paperwidth,ht=2.7ex,dp=1.2ex,rightskip=.5cm plus1fil]{date in head/foot}
      \color{upazul}\scriptsize\hfill\insertframenumber/\inserttotalframenumber
    \end{beamercolorbox}}
  \vskip0pt}

\newcommand{\concepto}[1]{\textcolor{upazul}{\textbf{#1}}}
\newcommand{\resultado}[1]{\textcolor{upverde}{\textbf{#1}}}
\newcommand{\advertencia}[1]{\textcolor{uprojo}{\textbf{#1}}}
\newcommand{\codigo}[1]{\colorbox{upazulclaro}{\strut\texttt{#1}}}
\newcommand{\fuente}[1]{\vfill\hfill{\tiny\color{upgris}#1}}


\usepackage{mathrsfs,bm,listings}
\lstset{language=Python,basicstyle=\ttfamily\footnotesize,columns=fullflexible,keepspaces=true,showstringspaces=false,breaklines=true}

\newcommand{\ovl}{\overline}
\newcommand{\ds}{\displaystyle}
\newcommand{\eps}{\varepsilon}

\DeclareMathOperator{\rd}{rd}
\DeclareMathOperator{\fl}{fl}
\DeclareMathOperator{\sign}{sign}
\DeclareMathOperator{\EPS}{eps}


\newcommand{\estirar}{\renewcommand{\arraystretch}{0.6}}

\newcommand{\Frac}{\displaystyle \frac}
\newcommand{\mfrac}[2]{\mbox{$\frac{#1}{#2}$}}
\newcommand{\ce}{\mathbb{C}}
\newcommand{\gment}[1]{\left\lceil #1 \right\rceil}

\newcommand{\rojo}[1]{\textcolor{uprojo}{#1}}

\title{Interpolación polinomial}
\subtitle{Lagrange, Newton, Hermite y aplicaciones}
\author{Oswaldo José Velásquez Castañón}
\institute{Universidad del Pacífico}
\date{Métodos Numéricos -- 2026-2}

\newtheorem{prop}{Proposici\'on}
\newtheorem{defi}{Definici\'on}

\newcommand{\tq}{\mid}
\newcommand{\re}{\mathbb{R}}
\newtheorem{thm}{Teorema}

\newtheorem{cor}{Corolario}

\theoremstyle{remark}


\newtheorem{ejem}{Ejemplo}


\DeclareMathOperator*{\res}{Res}
\newcommand{\sigmap}[2]{\sigma >_{#1} #2}
\newcommand{\zeros}{\partial}
\newcommand{\zeross}{\partial}
%\newcommand{\zerosa}{\mathfrak{Z}^*}
\newcommand{\zerosa}{\partial^*}
\newcommand{\zerossa}{\partial^*}
\newcommand{\K}{K}
\newcommand{\N}{\mathbb{N}}
\newcommand{\vct}[1]{\bm{#1}}
\newcommand{\Nr}{\mathbb{N}^{r*}}
\newcommand{\conv}{\circ}


\newcommand{\gamman}{\Gamma}

\DeclareMathOperator{\Mod}{mod}
\DeclareMathOperator{\Real}{Re}
\DeclareMathOperator{\Imag}{Im}
\renewcommand{\mod}{\Mod}


\DeclareMathOperator{\PSL}{PSL}
\DeclareMathOperator*{\Res}{Res}

\newcommand{\summ}{\displaystyle \sum}

\newcommand{\xframe}[1]{}


%\DeclareMathOperator{\rd}{rd}
%\DeclareMathOperator{\fl}{fl}
%\DeclareMathOperator{\sign}{sign}


\newcommand{\compacto}{\footnotesize
\setlength{\abovedisplayskip}{4pt}
\setlength{\belowdisplayskip}{4pt}
\setlength{\abovedisplayshortskip}{2pt}
\setlength{\belowdisplayshortskip}{2pt}}
\begin{document}

\begin{frame}[plain]
\begin{columns}[T,onlytextwidth]
\begin{column}{.70\textwidth}
\vspace{.7cm}{\usebeamerfont{title}\color{upazul}\inserttitle\par}
\vspace{.35cm}{\usebeamerfont{subtitle}\color{upgris}\insertsubtitle\par}
\vspace{1cm}{\large\insertauthor\par}{\normalsize\insertdate\par}
\end{column}
\begin{column}{.25\textwidth}\vspace{.55cm}\includegraphics[width=\linewidth]{UPlogo.pdf}\end{column}
\end{columns}
\end{frame}



%\section[Programme]{}
%\frame{\tableofcontents}

%\subsection{Overview of the Beamer Class}

\frame{
\frametitle{Interpolaci\'on polinomial}
Supongamos que tenemos $x_0,x_1,\dots,x_n$, $(n+1)$ puntos distintos del 
eje real y $ f_0,f_1,\dots,f_n$ los valores correspondientes de la ordenada. Buscamos un polinomio 
de grado menor o igual a $n$ %\index{Polinomio interpolante}
$$ P_n (x) = a_0+a_1 x+a_2 x^2 + \dots + a_n x^n $$
tal que $ P_n(x_i)=f_i, \quad i=0,1,\dots,n $.
\uncover<2>{
\begin{thm}
%\index{Polinomio interpolante!existencia y unicidad del}
Para $(n+1)$ puntos $(x_0,f_0), (x_1,f_1), \dots, (x_n,f_n) \in \re^2$ con 
abscisas distintas (esto es, $x_i\neq x_j$ para $i \neq j$), existe un \'unico polinomio 
$P_n $ de grado menor o igual a $ n $ que cumple las condiciones interpolantes:
$$ P_n(x_i)=f_i, \quad i=0,1,\dots,n $$
\end{thm}
}
}

\frame{
De hecho, las condiciones interpolantes
$$ a_0 + a_1 x_i + a_2 x_i^2 + \dots + a_n x_i^n = f_i, \quad i=0,1,\dots,n $$
constituyen un sistema lineal de $n+1$ ecuaciones con $n+1$ inc\'ognitas (los 
coeficientes del polinomio), que escribimos en la forma
$$ \left[
\begin{array}{ccccc}
1 & x_0 & x_0^2 & \dots &x_0^n \\
1 & x_1 & x_1^2 & \dots &x_1^n \\
\vdots & & & & \vdots \\
1 & x_n & x_n^2 & \dots &x_n^n \\
\end{array}
\right] \left[
\begin{array}{c}
a_0 \\ a_1 \\ \vdots \\ a_n 
\end{array}
\right] = \left[
\begin{array}{c}
f_0 \\ f_1 \\ \vdots \\ f_n 
\end{array}
\right].
$$
La matriz $(n+1) \times (n+1)$ del sistema es inversible, y por lo tanto el sistema posee soluci\'on \'unica.
}

\frame{
\compacto
Consideramos los siguientes datos:
$$ \begin{tabular}{c|ccccc}
$j$   &  $0$  &  $1$  &  $2$  & $3$ \\ \hline
$x_j$ &  $1$  & $3$   &  $5$  & $7$ \\
$f_j$ & $2.5$ & $3.5$ &  $7.5$  & $5.5$ 
\end{tabular} $$
\uncover<2->{El \'unico polinomio interpolante para esta tabla de valores tiene la forma
$$ P_3(x) = a_0 + a_1 x + a_2 x^2 + a_3 x^3. $$
}
\uncover<3->{Los coeficientes $a_0, a_1, a_2$ y $a_3$ resuelven el sistema
$$
\left[ \begin{array}{ccccc} 
1 & 1 & 1 & 1 \\
1 & 3 & 9 & 27\\
1 & 5 & 25 & 125 \\
1 & 7  & 49 & 343
\end{array}\right]
\left[
\begin{array}{c}
a_0 \\ a_1 \\ a_2 \\ a_3 
\end{array}
\right] = \left[
\begin{array}{c}
2.5 \\ 3.5 \\ 7.5 \\ 5.5
\end{array}
\right]
$$
}
\uncover<4>{Resolviendo el sistema obtenemos $a_0=5.93750$, $a_1=-5.31250$, $a_2=2.06250$ y $a_3=-0.18750$, de donde
$$ 
P_3(x)=5.93750-5.31250 x + 2.06250x^2 -0.18750 x^3.
$$
}
}


\frame{
\frametitle{F\'ormula de Lagrange}
Si $(x_k,f_k)$, $k=0,\dots,n$ son los $n+1$ puntos antes considerados, el polinomio interpolante se puede obtener directamente por la f\'ormula
$$ P_n (x) = f_0 L_0(x) +f_1 L_1(x) + \dots + f_n L_n(x) $$
donde los polinomios $ L_0,L_1,\dots,L_n $ tienen grado $ n $, 
%$$ f_j = P_n (x_j) = f_0 L_0(x_j) +f_1 L_1(x_j) + \dots + f_n L_n(x_j), \quad j=0,1,\dots,n $$
%Esto se cumple, 
si se exige que
$$ L_j(x_k) = \left\{
\begin{array}{cc}
0, & j \neq k \\
1, & j=k 
\end{array}
\right.
%= \delta_{jk} 
$$
Esto se cumple para los \emph{polinomios b\'asicos de Lagrange}
$$ L_j(x) = \prod_{ \substack{k=0 \\ k \neq j}}^n \frac{x-x_k}{x_j-x_k}. $$
%\tiny \begin{array}{c} k \neq j \\[-8ex] k=0 \end{array}

}


\frame{
\compacto
Consideramos los datos de la tabla
$$ \begin{tabular}{c|ccc}
$j$   & $0$ &  $1$  & $2$ \\ \hline
$x_j$ & $0$ & $0.5$ & $1$ \\ %\hline
$f_j$ & $1.1$ & $0.9$ & $0.6$
\end{tabular} $$
\uncover<2->{Los polinomios b\'asicos de Lagrange son
$$ \begin{array}{rcl}
L_0(x) &=& \Frac{(x-x_1)(x-x_2)}{(x_0-x_1)(x_0-x_2)} = \Frac{(x-0.5)(x-1)}{(0-0.5)(0-1)} = 2x^2-3x+1\\
L_1(x) &=& \Frac{(x-x_0)(x-x_2)}{(x_1-x_0)(x_1-x_2)} = \Frac{(x-0)(x-1)}{(0.5-0)(0.5-1)} = -4x^2+4x\\
L_2(x) &=& \Frac{(x-x_0)(x-x_1)}{(x_2-x_0)(x_2-x_1)} = \Frac{(x-0)(x-0.5)}{(1-0)(1-0.5)} = 2x^2-x
\end{array} $$
}
\uncover<3>{
Luego:
%$$ 
\begin{align*}
P_2(x) & = 1.1 (2x^2-3x+1) + 0.9 (-4x^2+4x) + 0.6 (2x^2-x) \\
	& = -0.2 x^2-0.3x+1.1
\end{align*}
}
}



\frame{
\compacto
\compacto
\frametitle{Forma de Newton}
Se inicia con el polinomio de grado $ 0 $ (constante) que pasa por el punto $ (x_0,f_0) $
$$ P_0 (x) = f_0. $$
A\~nadiendo $ (x_1,f_1) $, se construye un polinomio $ P_1(x) $ de grado no mayor que $ 1 $ tal que 
$$ P_1 (x_0) = f_0, \quad P_1 (x_1) = f_1. $$
Sea $ P_1(x) = P_0(x) + q_1(x), $ donde $grado(q_1) \leq 1 $. Como $ P_1 (x_0) = f_0 $, entonces 
$ q_1(x_0) = 0 $, luego $x-x_0$ divide $q_1(x)$ y 
$ q_1(x) = (x-x_0) \alpha(x)$.
Pero el grado de $q_1$ es $\leq 1$, luego $ \alpha(x) = \alpha $ (constante). Como $ P_1(x_1)=f_1 $, entonces 
$ f_1 = P_1(x_1)= f_0+ \alpha (x_1-x_0)$. De esto \'ultimo %$ \alpha $:
$$ \alpha = \Frac{f_1-f_0}{x_1-x_0} $$ %= \Frac{f_1-P_0(x_1)}{x_1-x_0}$$
Denotando $f[x_0]=f_0$, $f[x_1]=f_1$ y $\alpha=f[x_0,x_1]$, se tiene que
$$ f[x_0,x_1] = \frac{f[x_1]-f[x_0]}{x_1-x_0} $$
y
$$ P_1(x) = f[x_0]+f[x_0,x_1](x-x_0). $$
}

\frame{

\compacto
Ahora supongamos que tenemos $n+1$ puntos interpolantes 
$(x_i,f_i)$ con $0 \leq i \leq n$, y el polinomio $P_{n-1}(x)$ de grado $\leq n-1$
que interpola los $n$ primeros puntos, esto es
$$ P_{n-1}(x_i) = f_i, \quad 0 \leq i \leq n-1. $$
De manera an\'aloga a la anterior, se muestra que el polinomio $P_n(x)$ que interpola todos los puntos (incluido $(x_n,f_n)$) toma la forma

$$ P_n(x) = P_{n-1}(x) + f[x_0,x_1,\dots,x_n] (x-x_0)(x-x_1)\cdots(x-x_{n-1}), $$
donde $f[x_0,x_1,\dots,x_n]$ es el coeficiente de $x^n$ en el desarrollo de $P_n(x)$.
%De aqu\'i concluimos que la forma del polinomio interpolante es
%$$ P_n(x) = f[x_0] + \sum_{k=1}^{n-1} f[x_0,x_1,\dots,x_k] (x-x_0)(x-x_1)\cdots(x-x_{k-1}) $$
Reutilizando la recurrencia, el polinomio interpolante $ P_n $ puede ser escrito %entonces como:
como
%\index{Polinomio interpolante!de Newton}
$$ P_n(x) = f[x_0] + f[x_0,x_1] (x-x_0) + \dots + f[x_0,x_1,\dots,x_n] (x-x_0) \dots (x-x_{n-1}). $$


% $ P_n(x) = P_{n-1}(x) + q_{n}(x)$, donde $q_(x)$ es de grado $\leq n$. De las condiciones interpolantes
% $$ q_n(x_i) = P_n(x_i) - P_{n-1}(x_i) = f_i - f_i = 0, \mbox{ para } 0 \leq i \leq n-1. $$
% Luego cada factor $x-x_i$ divide a $q_n$ y $q(x) = \displaystyle \alpha(x) \prod_{k=0}^{n-1} (x-x_k)$; como antes analizando los grados de los polinomios, se concluye que $\alpha(x) = \alpha$ es constante.
% Si $f[x_0,x_1,\dots,x_n]$ es el coeficiente de $x^n$ en el desarrollo de $x^n$, entonces este coeficiente es igual a $\alpha$ ($P_{n-1}(x)$ no contiene un t\'ermino $x^n$).
% Conclusi\'on

}


\frame{
Denotemos con $ f[x_j,x_{j+1},\dots,x_{j+k}] $ al coeficiente de $x^k$ en el polinomio interpolante que pasa por los puntos 
$(x_j,f_j), (x_{j+1},f_{j+1}), \dots, (x_{j+k},f_{j+k})$. Se cumple:
$$ f[x_0] =f_0, \quad f[x_1] =f_1, \quad \dots , \quad f[x_n] =f_n $$
y se muestra que
%$$ f[x_0,\dots,x_n] = \Frac{f[x_1,\dots,x_n]-f[x_0,\dots,x_{n-1}]}{x_n-x_0} $$
%En general:
$$ f[x_j,\dots,x_{j+k}] = \Frac{f[x_{j+1},\dots,x_{j+k}]-f[x_j,\dots,x_{j+k-1}]}{x_{j+k}-x_j} $$

}

\frame{
Para calcular $ f[x_0,\dots,x_n] $, se utiliza la 
\emph{tabla de diferencias divididas}:
%\index{Diferencias divididas!tabla de}
\\[0.5cm]
\footnotesize{
\setlength{\tabcolsep}{2pt}
\begin{tabular}{cccccc} 
\hline 
 $x$    & $f[x]$          & $f[x,x]$     & $f[x,x,x]$       & $f[x,x,x,x]$         & $\dots$ \\
\hline 
$x_0$   & \rojo{$f_0 := f[x_0]$} \\
        &                 & \rojo{$f[x_0,x_1]=\frac{f[x_1]-f[x_0]}{x_1-x_0}$} \\
$x_1$   & $f_1 := f[x_1]$ &              & \rojo{$f[x_0,x_1,x_2]=\frac{f[x_1,x_2]-f[x_0,x_1]}{x_2-x_0}$} \\
        &                 & $f[x_1,x_2]=\frac{f[x_2]-f[x_1]}{x_2-x_1}$ &                  & \rojo{$f[x_0,x_1,x_2,x_3]$} \\
$x_2$   & $f_2 := f[x_2]$ &              & $f[x_1,x_2,x_3]=\frac{f[x_2,x_3]-f[x_1,x_2]}{x_3-x_1}$ &                      & $\dots$ \\
        &                 & $f[x_2,x_3]=\frac{f[x_3]-f[x_2]}{x_3-x_2}$ &                  & $f[x_1,x_2,x_3,x_4]$ \\
$x_3$   & $f_3 := f[x_3]$ &              & $f[x_2,x_3,x_4]=\frac{f[x_3,x_4]-f[x_2,x_3]}{x_4-x_2}$ & \vdots \\
        &                 & $f[x_3,x_4]=\frac{f[x_4]-f[x_3]}{x_4-x_3}$ & \vdots \\
$x_4$   & $f_4 := f[x_4]$ & \vdots \\ 
$\vdots$ & $\vdots$ \\
$x_n$   & $f_n := f[x_n]$ \\ 
\hline
\end{tabular}
}
}

\frame{
\compacto
\compacto
\renewcommand{\arraystretch}{0.9}
Retomamos el ejemplo de la tabla:
$$ \begin{tabular}{c|ccccc}
$j$   &  $0$  &  $1$  &  $2$  & $3$ \\ \hline
$x_j$ &  $1$  & $3$   &  $5$  & $7$ \\
$f_j$ & $2.5$ & $3.5$ &  $7.5$  & $5.5$ 
\end{tabular} $$
\uncover<2->{El esquema de diferencias divididas es
\begin{center}
\begin{tabular}{ccccc} 
\hline 
 $x$    & $f[x]$          & $f[x,x]$     & $f[x,x,x]$       & $f[x,x,x,x]$        \\
\hline 
$1$   & \rojo{$f[1]=2.5$} \\
       &                 & \rojo{$f[1,3]=0.5$} \\
$3$   & $f[3]=3.5$ &              & \rojo{$f[1,3,5]=0.375$} \\
        &                 & $f[3,5]=2$ &                  & \rojo{$f[1,3,5,7]=-0.1875$} \\
$5$   & $f[5]=7.5$ &              & $f[3,5,7]=-0.75$ \\ %&                      & $\dots$ \\
        &                 & $f[5,7]=-1$ \\ %&                  & $f[x_1,x_2,x_3,x_4]$ \\
$7$   & $f[7]=5.5$ \\
\hline
\end{tabular}
\end{center}
}
\uncover<3>{
El polinomio interpolante es
\begin{align*}
P_3(x) &= f[1] + f[1,3](x-1) + f[1,3,5](x-1)(x-3) \\
	& \qquad  + f[1,3,5,7](x-1)(x-3)(x-5)\\
	&= 2.5 + 0.5(x-1) +0.375 (x-1)(x-3) -0.1875(x-1)(x-3)(x-5)
\end{align*}
}
}


\frame{
\frametitle{Error de la aproximaci\'on polinomial}
Suponemos ahora que $f: \re \to \re$ es una funci\'on dada, y el polinomio interpolante 
$P_n(x)$ la interpola en $n+1$ puntos dados. Puesto que, en general 
$f(x) \neq P_n(x)$, hay un error en la aproximaci\'on de $f(x)$ mediante $P_n(x)$.
Como
$$f(x) = P_n(x) + f[x_0,x_1,\dots,x_n,x] \prod_{i=0}^n (x-x_i), $$
el error en la aproximaci\'on por el polinomio $P_n(x)$ (y que pasa por los $n+1$ puntos prefijados) 
est\'a dado por
$$ E(x) = f[x_0,x_1,\dots,x_n,x] \prod_{i=0}^n (x-x_i). $$
\rojo{Podemos calcular este error empleando un punto adicional $x$ en el esquema 
de diferencias divididas, si conocemos $f(x)$.}
}



\begin{frame}{Teorema del error de interpolación}
\begin{thm}
Sea $f\in C^{n+1}[a,b]$, con nodos distintos $x_0,\dots,x_n\in[a,b]$.
Sea $P_n$ el polinomio de grado $\leq n$ que satisface
$P_n(x_i)=f(x_i)$, $i=0,\dots,n$.
Para cada $x\in[a,b]$ que no sea un nodo, existe
\[
\xi\in\bigl(\min\{x,x_0,\dots,x_n\},
                 \max\{x,x_0,\dots,x_n\}\bigr)
\]
tal que
\[
f(x)-P_n(x)=\frac{f^{(n+1)}(\xi)}{(n+1)!}
                  \prod_{i=0}^n(x-x_i).
\]
En los nodos, el error es cero.
\end{thm}
\end{frame}

\begin{frame}{Cota del error de interpolación}
Si $M_{n+1}=\max_{t\in[a,b]}|f^{(n+1)}(t)|$, entonces
\[
|f(x)-P_n(x)|\leq\frac{M_{n+1}}{(n+1)!}
                       \left|\prod_{i=0}^n(x-x_i)\right|.
\]
La cota depende tanto de la función como de la distribución de nodos.
\begin{block}{Error de aproximación y error de cálculo}
La fórmula supone que evaluamos el interpolante exacto.
El cálculo en punto flotante agrega errores de redondeo.
\end{block}
El residual en los nodos verifica las condiciones interpolantes,
pero no mide el error entre los nodos.
\end{frame}

\frame{
Nota sobre interpolaci\'on con condiciones sobre las derivadas:

Observe que 
$$ \lim_{h \to 0} \frac{f[x+h]-f[x]}{h} = f'(x), $$
por lo que definimos
$$ f[x,x] = f'(x). $$
Igualmente
$$ f[x,x,x] = \frac{1}{2} f''(x), \quad f[x,x,x,x] = \frac{1}{6} f'''(x) $$
y, en general (para $n+1$ puntos iguales $x$)
$$ f[x,x,\dots, x] = \frac{1}{n!} f^{(n)}(x). $$
Esto nos permite reemplazar condiciones sobre las derivadas en la tabla de diferencias divididas, 
siempre y cuando el n\'umero de condiciones totales sea igual al n\'umero de coeficientes a calcular. La interpolaci\'on con este tipo de condiciones se denomina \emph{interpolaci\'on de Hermite}.
}

\frame{
Consideremos la tabla
\begin{center}
\begin{tabular}{c|cc}
$x_i$ & $x_0$ & $x_1$ \\ \hline
$f_i$ & $f_0$ & $f_1$ \\ \hline
$f'_i$ & $f'_0$ & $f'_1$ \\ \hline
$f''_i$ & $f''_0$ & $f''_1$ \\
\end{tabular}
\end{center}
Tenemos $ 6$ condiciones interpolantes. Se est\'a buscando un polinomio de 
grado $\leq 5$: $ P_5(x) $.
}


\frame{
%{ }\\[0.5cm]
\begin{center}
\tiny{\begin{tabular}{ccccccc} 
\hline 
 $x$    & $f[x]$   & $f[x,x]$     & $f[x,x,x]$        & $f[x,x,x,x]$         & $f[x,x,x,x,x]$           & $f[x,x,x,x,x,x]$ \\
\hline 
$x_0$   & $f[x_0]$ \\
        &          & $f'(x_0) $ \\
$x_0$   & $f[x_0]$ &              & $\Frac{1}{2}f''(x_0)$ \\
        &          & $f'(x_0)   $ &                   & $f[x_0,x_0,x_0,x_1]$ \\
$x_0$   & $f[x_0]$ &              & $f[x_0,x_0,x_1]$ &                      & $f[x_0,x_0,x_0,x_1,x_1]$ \\
        &          & $f[x_0,x_1]$ &                   & $f[x_0,x_0,x_1,x_1]$ &                          & $f[x_0,x_0,x_0,x_1,x_1,x_1]$ \\
$x_1$   & $f[x_1]$ &              & $f[x_0,x_1,x_1]$ &                      & $f[x_0,x_0,x_1,x_1,x_1]$ \\
        &          & $f'(x_1)   $ &                   & $f[x_0,x_1,x_1,x_1]$ \\
$x_1$   & $f[x_1]$ &              & $\Frac{1}{2}f''(x_1)$ \\
        &          & $f'(x_1)   $ \\
$x_1$   & $f[x_1]$  \\
\hline
\end{tabular}}
\end{center}
Entonces 
\begin{align*} 
P_5(x) =& b_0 + b_1(x-x_0) + b_2(x-x_0)^2 + b_3(x-x_0)^3 \\
	&+ b_4(x-x_0)^3 (x-x_1) + b_5 (x-x_0)^3 (x-x_1)^2 
\end{align*}
%$$
donde los $ b_i, \quad i=0,\dots,5 $ son los valores de la primera l\'{\i}nea oblicua.
}


\frame{
\frametitle{Spline c\'ubico}

Sean los puntos $(x_i,f_i)$, $i=1,\dots,n$ con 
$$a=x_1 < x_2 < \dots < x_{n-1} < x_n = b. $$
Un \emph{spline c\'ubico} es una funci\'on $S:[a,b] \to \re$ tal que
\begin{enumerate}
 \item $S(x)$ es un polinomio de grado $\leq 3$ en cada intervalo $[x_j,x_{j+1}]$, 
$j=1,2,\dots,n-1$;
\item $S(x_i)=f_i$, $i=1,2,\dots,n$;
\item $S$, $S'$ y $S''$ son funciones continuas en $[a,b]$ (el spline $S$ es dos veces continuamente diferenciable en $[a,b]$).
\end{enumerate}

}

\frame{
\frametitle{Tipos de splines}
Un spline c\'ubico es
\begin{enumerate}
 \item \emph{natural} si
$$ S''(a) = S''(b) = 0, $$
\item \emph{anclado} si 
$$ S'(a) = \alpha, \quad S'(b) = \beta $$
para ciertos valores $\alpha, \beta$.
\item \emph{peri\'odico} si $S$ es una funci\'on peri\'odica de periodo $b-a$, esto es
$$ S(a)=S(b), \quad S'(a)=S'(b), \quad S''(a)=S''(b). $$
\end{enumerate}
}

\frame{
Determine si la siguiente funci\'on es un spline c\'ubico:
$$
f(x) = \left\{
\begin{array}{ll}
x^3+x, & x \in [-1,0] \\
x^3-x, & x \in [0,1].
\end{array}
\right.
$$
}

\frame{
Determine todos los valores de $a$, $b$, $c$, $d$ y $e$ para los cuales la siguiente funci\'on es un spline c\'ubico:
$$
f(x) = \left\{
\begin{array}{ll}
a(x-2)^2 + b(x-1)^3, & x \in [-10, 1] \\
c(x-2)^2, & x \in [1,3] \\
d(x-2)^2 + e(x-3)^3, & x \in [3,10]. 
\end{array}
\right.
$$
}


\frame{
\compacto
\compacto
\frametitle{Construcci\'on del spline c\'ubico}

Observe que si un spline $S$ tiene $n$ nodos, estar\'a definido por $n-1$ polinomios c\'ubicos, de donde se necesitan $4(n-1)$ coeficientes. %Un spline interpolante para $5$ puntos necesita del c\'alculo de $16$ coeficientes.

Definimos $h_i=x_{i+1}-x_i$ para $i=1,2,\dots,n-1$. 
En cada intervalo $[x_i,x_{i+1}]$, desarrollamos el esquema de diferencias divididas siguiente:
\begin{center}
\begin{tabular}{ccccc} 
%\hline 
% $x$    & $f[x]$          & $f[x,x]$     & $f[x,x,x]$       & $f[x,x,x,x]$        \\
\hline 
$x_i$     & $f_i$ \\
          &                 & \rojo{$\lambda_i=?$} \\
$x_i$     & $f_i$ &              & $\frac{m_i-\lambda_i}{h_i}$ \\
          &                 & $m_i$ &                  & $\frac{\lambda_{i+1}-2m_i+\lambda_i}{h_i^2}$ \\
$x_{i+1}$ & $f_{i+1}$ &                    & $\frac{\lambda_{i+1}-m_i}{h_i}$ \\ %&                      & $\dots$ \\
          &                 & $\lambda_{i+1}$ \\ %&                  & $f[x_1,x_2,x_3,x_4]$ \\
$x_{i+1}$ & $f_{i+1}$ \\
\hline
\end{tabular}
\end{center}
El polinomio c\'ubico que representa el spline en $[x_i,x_{i+1}]$ est\'a dado por
$$ S(x) = f_i + \lambda_i(x-x_i) + \frac{m_i-\lambda_i}{h_i} (x-x_i)^2 + 
\frac{\lambda_{i+1}-2m_i + \lambda_i}{h_i^2} (x-x_i)^2 (x-x_{i+1}). $$

}

\frame{
\compacto
Para un spline natural, con $m_i=(f_{i+1}-f_i)/h_i$ y $\lambda_i=S'(x_i)$, debemos resolver el sistema lineal
$$
\begin{array}{l}
2 \lambda_1 + \lambda_2 = 3m_1 \\
\frac{h_i}{h_{i-1}} \lambda_{i-1} + 2 \Bigl( 1 + \frac{h_i}{h_{i-1}} \Bigr) \lambda_i + \lambda_{i+1} = 3 \Bigl( m_i + \frac{h_i}{h_{i-1}} m_{i-1} \Bigr), \quad i=2,\dots,n-1 \\
\lambda_{n-1} + 2 \lambda_n = 3 m_{n-1}
\end{array}
$$
}





\frame{
\frametitle{Diferenciaci\'on num\'erica} %\index{Diferenciaci\'on num\'erica}

Consideramos una funci\'on $f$ interpolada  por un polinomio $P_n(x)$ en $n+1$ puntos.
Sean $ (x_i,f_i), \quad i=0,1,\dots,n $ puntos interpolantes. De la f\'ormula de Lagrange:
$$ P_n(x) = \sum_{j=0}^n f_j \alpha_j \prod_{\substack{k=0 \\ k \neq j}}^n (x-x_k) $$
con $ \alpha_j = \ds\prod_{\substack{k=0 \\ k \neq j}}^n \Frac{1}{x_j-x_k} $. \\
Se obtiene:
$$ \frac{d^n P_n}{dx^n}(x) = \sum_{j=0}^n f_j \alpha_j n! \sim f^{(n)}(x) $$
que resulta una aproximaci\'on de la n-\'esima derivada $ f^{(n)} $.
}

\frame{
\frametitle{F\'ormulas de diferencias para las derivadas:}
Consideramos puntos igualmente espaciados, con $x_{j+1}-x_j = h$. Entonces 
$ \alpha_j = \frac{(-1)^{n-j}}{h^n n!} \binom{n}{j}$, y obtenemos la f\'ormula
%y luego
$$ f^{(n)} (x_0) \approx \sum_{j=0}^n f_j \frac{(-1)^{n-j}}{h^n} \binom{n}{j}. $$
Para $ n=1,2,3$, con $f_j=f(x+jh)$:
$$ \begin{array}{l}
f'(x) \sim \frac{f_1-f_0}{h} = \frac{f(x+h)-f(x)}{h}\\[0.3cm]
f''(x) \sim \frac{f_2-2 f_1+f_0}{h^2} = \frac{f(x+2h)-2f(x+h)+f(x)}{h^2}\\[0.3cm]
f'''(x) \sim \frac{f_3-3 f_2 + 3 f_1 - f_0}{h^3} = \frac{f(x+3h)-3f(x+2h)+3f(x+h)-f(x)}{h^3}
\end{array} $$

}



\frame{

\frametitle{Interpolaci\'on en dimensi\'on dos}

Dado un rect\'angulo $[a,b] \times [c,d]$, consideremos una malla regular, formada por puntos igualmente espaciados
en cada dimensi\'on
$$x_i= a+ i h, i=0,\dots,n, \quad y_j=c+jk, j=0,\cdots,m$$
con $h=\frac{b-a}{n}$, $k=\frac{d-c}{m}$. 

Denotamos
$$ p_{i,j} =(x_i,y_j), \quad f_{i,j} = f(p_{i,j}); $$
consideramos entonces el problema de interpolaci\'on en cada rect\'angulo
$$ R_{i,j} = [x_i,x_{i+1}] \times [y_j,y_{j+1}], \quad i=0,\dots,n-1, \; j=0,\dots,m-1. $$


}

\frame{
\compacto
La funci\'on 
%$$ 
\begin{align*}
P_{i,j}(x,y) =& f_{i,j} +mx_{i,j}(x-x_i)+my_{i,j} (y-y_j) \\
	&+ \frac{mx_{i,j+1}-mx_{i,j}}{k}(x-x_i)(y-y_j), 
\end{align*}
%$$
para $(x,y) \in R_{ij}$, donde 
$$ mx_{i,j}= \frac{f_{i+1,j}-f_{i,j}}{h}, \quad my_{i,j} = \frac{f_{i,j+1}-f_{i,j}}{k} $$
cumple las condiciones interpolantes en el rect\'angulo $R_{i,j}$. \\

Con las condiciones establecidas, la funci\'on $P:[a,b] \times [c,d] \to \re$ tal que 
$$ P(x,y) = P_{i,j}(x,y), \mbox{ para } (x,y) \in R_{i,j}, \quad i=0,\dots,n-1, \; j=0,\dots,m-1 $$
cumple las condiciones interpolantes, y es una funci\'on continua (en general, no diferenciable en las interfaces).

}

\frame{

\compacto
En cada rect\'angulo $R_{i,j}$, la funci\'on $P_{i,j}$ tambi\'en puede 
escribirse en forma similar a la forma de Lagrange 

$$ P(x,y) = f_{i,j} \alpha_{i,j}(x,y) + f_{i+1,j} \beta_{i,j}(x,y) 
+ f_{i+1,j+1} \gamma_{i,j}(x,y) + f_{i,j+1} \delta_{i,j}(x,y), $$
donde 
$$ \alpha_{i,j}(x,y) = \Bigl(1-\frac{x-x_i}{h} \Bigr) \Bigl(  1 - \frac{y-y_j}{k} \Bigr), $$
$$ \beta_{i,j}(x,y) = \Bigl( \frac{x-x_i}{h} \Bigr) \Bigl(  1 - \frac{y-y_j}{k} \Bigr), $$
$$ \gamma_{i,j}(x,y) = \Bigl( \frac{x-x_i}{h} \Bigr) \Bigl(  \frac{y-y_j}{k} \Bigr), $$
$$ \delta_{i,j}(x,y) = \Bigl( 1 -\frac{x-x_i}{h} \Bigr) \Bigl(   \frac{y-y_j}{k} \Bigr). $$
}

\frame{
\compacto
\frametitle{Interpolaci\'on en una malla triangulada}

Consideramos una malla triangular de puntos $p_i=(x_i,y_i)$, $i=1,\dots,n$. 
Dados tres puntos $p_i$, $p_j$, $p_k$ no colineales formando un tri\'angulo, formamos una funci\'on interpolante
$P$ sobre el tri\'angulo $\Delta = p_ip_jp_k$ 
$$ P(x,y) = f_i \phi_i(x,y) + f_j \phi_j(x,y) + f_k \phi_k(x,y), $$
donde 
$$ \phi_i(x,y)= F_{j,k}(x,y), \quad
\phi_j(x,y) = F_{k,i}(x,y), \quad
\phi_k(x,y) = F_{i,j}(x,y), $$
donde 
$$ F_{s,t} = \frac{1}{A_{i,j,k}} ( \alpha_{s,t}  + \beta_{s,t} x - \gamma_{s,t} y), $$
con 
$$ \alpha_{s,t} = x_s y_t - x_t y_s, \quad \beta_{s,t} = y_s - y_t, \quad \gamma_{s,t} = x_s - x_t, $$
y
$$ A_{i,j,k} = \left|
\begin{array}{ccc}
1 & x_i & y_i \\ 1 & x_j & y_j \\ 1 & x_k & y_k
\end{array}
\right|.$$


}


\frame{
\rojo{Observaci\'on. } De hecho, la gr\'afica de $P$ es, dentro del tri\'angulo 
$ \Delta = p_ip_jp_k$, el plano que pasa por los puntos $(x_i,y_i,f_i)$, 
$(x_j,y_j,f_j)$ y $(x_k,y_k,f_k)$. Podemos determinar la ecuaci\'on del plano haciendo
$$ P(x,y) = a + bx + cy, $$
y luego estableciendo las ecuaciones (condiciones interpolantes)
$$P(x_s,y_s) = a + bx_s + c y_s = f_s, \quad s=i,j,k $$
esto es 
$$ 
\left[
\begin{array}{ccc}
1 & x_i & y_i \\
1 & x_j & y_j \\
1 & x_k & y_k
\end{array}
\right]
\left[
\begin{array}{c}
 a \\ b \\ c 
\end{array}
\right] = 
\left[
\begin{array}{c}
 f_i \\ f_j \\ f_k 
\end{array}
\right] 
$$
para determinar $a,b,c$. Las f\'ormulas anteriores se obtienen aplicando la regla de Cramer para resolver este sistema.
}


\begin{frame}[fragile]{Vandermonde en NumPy}
Retomamos los mismos cuatro puntos de la presentación.
\begin{lstlisting}
import numpy as np
from numpy.polynomial.polynomial import polyval

x = np.array([1., 3., 5., 7.])
f = np.array([2.5, 3.5, 7.5, 5.5])
V = np.vander(x, increasing=True)
a = np.linalg.solve(V, f)
print(a)
# [5.9375, -5.3125, 2.0625, -0.1875]
print(polyval(4., a))  # 5.6875
\end{lstlisting}
\texttt{increasing=True} ordena las columnas como $1,x,x^2,x^3$.
\texttt{polyval} usa los coeficientes en ese mismo orden.
\end{frame}

\begin{frame}[fragile]{Diferencias divididas en NumPy}
La recurrencia anterior permite obtener los coeficientes de Newton.
\begin{lstlisting}
def diferencias_divididas(x, f):
    x = np.asarray(x, dtype=float)
    c = np.array(f, dtype=float, copy=True)
    if (x.ndim != 1 or c.ndim != 1
            or len(x) != len(c) or len(x) == 0):
        raise ValueError("Datos incompatibles")
    if len(np.unique(x)) != len(x):
        raise ValueError("Nodos repetidos")
    for k in range(1, len(x)):
        c[k:] = ((c[k:] - c[k-1:-1])
                 / (x[k:] - x[:-k]))
    return c
\end{lstlisting}
La función requiere nodos distintos. Hermite necesita un tratamiento
específico de los nodos repetidos.
\end{frame}

\begin{frame}[fragile]{Evaluación anidada de Newton}
Para cuatro nodos,
\[
P_3(t)=c_0+(t-x_0)\bigl[c_1+(t-x_1)
                  [c_2+(t-x_2)c_3]\bigr].
\]
\begin{lstlisting}
def evaluar_newton(t, x, c):
    t = np.asarray(t, dtype=float)
    p = np.full_like(t, c[-1], dtype=float)
    for k in range(len(c)-2, -1, -1):
        p = c[k] + (t-x[k])*p
    return p

c = diferencias_divididas(x, f)
print(c)  # [2.5, 0.5, 0.375, -0.1875]
print(evaluar_newton(4., x, c))  # 5.6875
\end{lstlisting}
Construcción: $O(n^2)$. Evaluación por punto: $O(n)$.
\end{frame}

\begin{frame}[fragile]{Verificación de la implementación}
\begin{lstlisting}
q = np.linspace(x.min(), x.max(), 101)
p_v = polyval(q, a)
p_n = evaluar_newton(q, x, c)

assert np.allclose(polyval(x, a), f)
assert np.allclose(evaluar_newton(x, x, c), f)
assert np.allclose(p_v, p_n)

residual = np.max(np.abs(polyval(x, a)-f))
condicion = np.linalg.cond(V)
print(residual, condicion)
\end{lstlisting}
Vandermonde y Newton representan el mismo polinomio en aritmética exacta.
Una matriz de Vandermonde mal condicionada puede amplificar perturbaciones.
La coincidencia de los valores en los nodos no garantiza una buena
aproximación de la función en todo el intervalo.
\end{frame}

\begin{frame}{Convenciones de NumPy}
\begin{itemize}
\item \texttt{numpy.polynomial.polynomial.polyval(t, a)} recibe
$a_0,a_1,\dots,a_n$ en orden creciente de potencias.
\item \texttt{numpy.polyval(a, t)} recibe los coeficientes en orden decreciente.
\item \texttt{numpy.interp} realiza interpolación lineal por tramos.
\item Un ajuste de grado menor que el número de nodos menos uno
resuelve un problema diferente de la interpolación exacta.
\end{itemize}
\vfill
{\scriptsize Documentación oficial de NumPy:
\href{https://numpy.org/doc/stable/reference/generated/numpy.vander.html}{vander},
\href{https://numpy.org/doc/stable/reference/generated/numpy.linalg.solve.html}{linalg.solve},
\href{https://numpy.org/doc/stable/reference/generated/numpy.polynomial.polynomial.polyval.html}{polynomial.polyval}.}
\end{frame}

\end{document}
